\documentclass[12pt,a4paper]{article}
\usepackage[utf8]{inputenc}
\usepackage[T1]{fontenc}
\usepackage{amsmath,amsfonts,amssymb}
\usepackage{graphicx}
\usepackage{booktabs}
\usepackage{longtable}
\usepackage{array}
\usepackage{multirow}
\usepackage{wrapfig}
\usepackage{float}
\usepackage{colortbl}
\usepackage{pdflscape}
\usepackage{tabu}
\usepackage{threeparttable}
\usepackage{threeparttablex}
\usepackage{forloop}
\usepackage{datetime}
\usepackage{hyperref}
\usepackage{geometry}
\usepackage{fancyhdr}
\usepackage{titlesec}
\usepackage{listings}
\usepackage{xcolor}
\usepackage{enumitem}
\usepackage{tikz}
\usepackage{pgfplots}
\usepackage{pgfplotstable}

% Page setup
\geometry{margin=1in}
\pagestyle{fancy}
\fancyhf{}
\rhead{DeFiMon Data Analytics Guide}
\lhead{Multi-Source ML Analysis}
\rfoot{Page \thepage}

% Title formatting
\titleformat{\section}{\Large\bfseries}{\thesection}{1em}{}
\titleformat{\subsection}{\large\bfseries}{\thesubsection}{1em}{}

% Code listing setup
\lstset{
    language=Python,
    basicstyle=\ttfamily\small,
    keywordstyle=\color{blue},
    commentstyle=\color{green!60!black},
    stringstyle=\color{red},
    numbers=left,
    numberstyle=\tiny,
    frame=single,
    breaklines=true,
    showstringspaces=false
}

% Custom colors
\definecolor{blockchain}{RGB}{52, 152, 219}
\definecolor{github}{RGB}{36, 41, 46}
\definecolor{social}{RGB}{29, 161, 242}
\definecolor{news}{RGB}{231, 76, 60}
\definecolor{defi}{RGB}{46, 204, 113}

\begin{document}

\title{\Huge\textbf{Comprehensive Data Analytics Guide for Machine Learning}\\
\Large Multi-Source Blockchain Data Analysis Framework}
\author{DeFiMon Analytics Team}
\date{\today}

\maketitle

\begin{abstract}
This document provides a comprehensive framework for data analysts working with multi-source blockchain data for machine learning applications. We present detailed methodologies for data collection, preprocessing, feature engineering, and analysis across five primary data sources: blockchain transactions, DeFi protocol metrics, GitHub development activity, social media sentiment, and news media coverage. The guide includes statistical analysis techniques, visualization methods, and machine learning preparation workflows specifically designed for cryptocurrency and DeFi analytics.
\end{abstract}

\tableofcontents
\newpage

\section{Introduction to Multi-Source Data Analytics}

\subsection{Data Ecosystem Overview}

Our multi-source data analytics framework integrates five primary data streams to provide comprehensive insights for machine learning applications in the blockchain and DeFi space. The system processes approximately 94GB of raw data daily, generating 2GB of ML-ready datasets with 125 engineered features.

\begin{table}[h]
\centering
\caption{Data Source Overview}
\begin{tabular}{|l|c|c|c|c|}
\hline
\textbf{Data Source} & \textbf{Volume (GB/day)} & \textbf{Update Frequency} & \textbf{Features} & \textbf{Quality Score} \\
\hline
Blockchain & 50 & Real-time & 50 & 98.5\% \\
DeFi Protocols & 15 & 5-min & 25 & 96.2\% \\
GitHub & 2 & 1-hour & 20 & 95.2\% \\
Social Media & 23 & Real-time & 30 & 87.3\% \\
News Media & 4 & 15-min & 25 & 92.1\% \\
\hline
\textbf{Total} & \textbf{94} & \textbf{Various} & \textbf{150} & \textbf{93.3\%} \\
\hline
\end{tabular}
\end{table}

\subsection{Analytical Framework Architecture}

The analytical framework consists of four primary stages:

\begin{enumerate}
    \item \textbf{Data Ingestion}: Real-time collection from multiple APIs and sources
    \item \textbf{Data Preprocessing}: Cleaning, validation, and standardization
    \item \textbf{Feature Engineering}: Creation of 125 ML-ready features
    \item \textbf{Model Preparation}: Dataset construction for machine learning
\end{enumerate}

\section{Blockchain Data Analysis}

\subsection{Data Structure and Schema}

Blockchain data constitutes our primary data source, providing real-time transaction and network metrics across 17 blockchain networks.

\subsubsection{Transaction Data Schema}

\begin{lstlisting}
{
    "transaction": {
        "hash": "0xabc123...",
        "block_number": 12345678,
        "from": "0x742d35cc6634c0532925a3b8d4c9db96c4b4d8b6",
        "to": "0x742d35cc6634c0532925a3b8d4c9db96c4b4d8b6",
        "value": "0x0",
        "gas_price": "0x59682f00",
        "gas_used": "0x5208",
        "gas_limit": "0x186a0",
        "nonce": "0x0",
        "input": "0x...",
        "timestamp": "2024-01-15T10:30:00Z"
    }
}
\end{lstlisting}

\subsubsection{Network Metrics Schema}

\begin{lstlisting}
{
    "network_metrics": {
        "block_number": 12345678,
        "timestamp": "2024-01-15T10:30:00Z",
        "gas_used": "0x123456",
        "gas_limit": "0x1c9c380",
        "difficulty": "0x123456789",
        "total_difficulty": "0x123456789abcdef",
        "size": "0x1234",
        "uncles": [],
        "base_fee_per_gas": "0x59682f00",
        "extra_data": "0x..."
    }
}
\end{lstlisting}

\subsection{Statistical Analysis Methods}

\subsubsection{Transaction Volume Analysis}

The transaction volume analysis examines patterns in blockchain activity to identify market trends and network congestion.

\begin{equation}
\text{Volume}_{24h} = \sum_{i=1}^{n} \text{Value}_i \times \text{Price}_i
\end{equation}

\begin{equation}
\text{Volume Change}_{24h} = \frac{\text{Volume}_{t} - \text{Volume}_{t-24h}}{\text{Volume}_{t-24h}} \times 100\%
\end{equation}

\subsubsection{Gas Price Analysis}

Gas price analysis provides insights into network congestion and transaction prioritization.

\begin{equation}
\text{Gas Price}_{avg} = \frac{1}{n} \sum_{i=1}^{n} \text{Gas Price}_i
\end{equation}

\begin{equation}
\text{Gas Price}_{percentile} = P_{95}(\text{Gas Price}_i)
\end{equation}

\subsubsection{Address Activity Analysis}

Address activity analysis tracks user behavior and network participation.

\begin{equation}
\text{Active Addresses}_{24h} = |\text{Unique}(\text{From Addresses} \cup \text{To Addresses})|
\end{equation}

\begin{equation}
\text{Transaction Frequency} = \frac{\text{Total Transactions}_{24h}}{\text{Active Addresses}_{24h}}
\end{equation}

\subsection{Feature Engineering for Blockchain Data}

\subsubsection{Price-Based Features}

\begin{lstlisting}
def calculate_price_features(data):
    """Calculate price-based features"""
    features = {}
    
    # Price changes
    features['price_change_1h'] = data['price'].pct_change(periods=1)
    features['price_change_24h'] = data['price'].pct_change(periods=24)
    features['price_change_7d'] = data['price'].pct_change(periods=168)
    
    # Price volatility
    features['price_volatility_1h'] = data['price'].rolling(1).std()
    features['price_volatility_24h'] = data['price'].rolling(24).std()
    features['price_volatility_7d'] = data['price'].rolling(168).std()
    
    # Price momentum
    features['price_momentum'] = data['price'] / data['price'].shift(24) - 1
    
    return features
\end{lstlisting}

\subsubsection{Volume-Based Features}

\begin{lstlisting}
def calculate_volume_features(data):
    """Calculate volume-based features"""
    features = {}
    
    # Volume metrics
    features['volume_24h'] = data['volume'].rolling(24).sum()
    features['volume_change_24h'] = features['volume_24h'].pct_change()
    features['volume_sma_7d'] = data['volume'].rolling(168).mean()
    
    # Volume-price relationship
    features['volume_price_correlation'] = (
        data['volume'].rolling(24)
        .corr(data['price'])
    )
    
    # Volume volatility
    features['volume_volatility'] = data['volume'].rolling(24).std()
    
    return features
\end{lstlisting}

\subsubsection{Network-Based Features}

\begin{lstlisting}
def calculate_network_features(data):
    """Calculate network-based features"""
    features = {}
    
    # Gas metrics
    features['gas_price_avg'] = data['gas_price'].rolling(24).mean()
    features['gas_price_median'] = data['gas_price'].rolling(24).median()
    features['gas_price_std'] = data['gas_price'].rolling(24).std()
    
    # Network congestion
    features['gas_used_ratio'] = data['gas_used'] / data['gas_limit']
    features['pending_transactions'] = data['pending_tx_count']
    features['mempool_size'] = data['mempool_size']
    
    # Block metrics
    features['block_time_avg'] = data['block_time'].rolling(24).mean()
    features['block_time_std'] = data['block_time'].rolling(24).std()
    features['uncle_rate'] = data['uncle_blocks'] / data['total_blocks']
    
    return features
\end{lstlisting}

\section{DeFi Protocol Data Analysis}

\subsection{Protocol Metrics Schema}

DeFi protocol data provides insights into decentralized finance activity across multiple protocols.

\begin{lstlisting}
{
    "protocol_metrics": {
        "protocol": "uniswap-v3",
        "timestamp": "2024-01-15T10:30:00Z",
        "tvl": 1234567.89,
        "volume_24h": 987654.32,
        "fees_24h": 12345.67,
        "users_24h": 5432,
        "transactions_24h": 12345,
        "pools_count": 1234,
        "tokens_count": 567
    }
}
\end{lstlisting}

\subsection{TVL Analysis}

Total Value Locked (TVL) analysis examines the capital efficiency and growth of DeFi protocols.

\begin{equation}
\text{TVL Change}_{24h} = \frac{\text{TVL}_{t} - \text{TVL}_{t-24h}}{\text{TVL}_{t-24h}} \times 100\%
\end{equation}

\begin{equation}
\text{TVL Growth Rate} = \frac{\ln(\text{TVL}_{t}) - \ln(\text{TVL}_{t-7d})}{7}
\end{equation}

\subsection{Yield Analysis}

Yield analysis examines the returns generated by different DeFi protocols.

\begin{equation}
\text{APY}_{protocol} = \frac{\text{Annual Revenue}}{\text{TVL}} \times 100\%
\end{equation}

\begin{equation}
\text{Yield Spread} = \text{APY}_{protocol} - \text{Risk Free Rate}
\end{equation}

\subsection{Liquidity Analysis}

Liquidity analysis examines the depth and efficiency of trading pools.

\begin{equation}
\text{Liquidity Depth} = \sum_{i=1}^{n} \text{Reserve}_i \times \text{Price}_i
\end{equation}

\begin{equation}
\text{Impermanent Loss} = \frac{\text{Value}_{HODL} - \text{Value}_{LP}}{\text{Value}_{HODL}} \times 100\%
\end{equation}

\section{GitHub Development Activity Analysis}

\subsection{Repository Metrics Schema}

GitHub data provides insights into development activity and project health.

\begin{lstlisting}
{
    "repository_metrics": {
        "repo_id": 123456789,
        "name": "uniswap-v3-core",
        "full_name": "Uniswap/uniswap-v3-core",
        "language": "Solidity",
        "stars": 5000,
        "forks": 1200,
        "watchers": 500,
        "open_issues": 45,
        "open_pulls": 23,
        "last_commit": "2024-01-15T10:30:00Z",
        "created_at": "2021-03-23T00:00:00Z",
        "updated_at": "2024-01-15T10:30:00Z"
    }
}
\end{lstlisting}

\subsection{Development Velocity Analysis}

Development velocity measures the pace and quality of code development.

\begin{equation}
\text{Commit Frequency} = \frac{\text{Commits}_{7d}}{7}
\end{equation}

\begin{equation}
\text{Code Churn Rate} = \frac{\text{Additions} + \text{Deletions}}{\text{Total Lines}} \times 100\%
\end{equation}

\begin{equation}
\text{Issue Velocity} = \frac{\text{Issues Closed}_{7d}}{\text{Issues Opened}_{7d}}
\end{equation}

\subsection{Community Health Analysis}

Community health analysis examines the engagement and sustainability of open-source projects.

\begin{equation}
\text{Star Growth Rate} = \frac{\text{Stars}_{t} - \text{Stars}_{t-7d}}{7}
\end{equation}

\begin{equation}
\text{Contributor Diversity} = \frac{\text{Unique Contributors}_{30d}}{\text{Total Commits}_{30d}}
\end{equation}

\begin{equation}
\text{Maintenance Score} = \frac{\text{Issues Closed}_{30d}}{\text{Issues Opened}_{30d}} \times \text{PR Merge Rate}
\end{equation}

\subsection{Code Quality Analysis}

Code quality analysis examines the technical health of codebases.

\begin{equation}
\text{Bug Ratio} = \frac{\text{Bug Issues}}{\text{Total Issues}} \times 100\%
\end{equation}

\begin{equation}
\text{Test Coverage} = \frac{\text{Tested Lines}}{\text{Total Lines}} \times 100\%
\end{equation}

\begin{equation}
\text{Security Score} = \frac{\text{Security Audits}}{\text{Total Releases}} \times 100\%
\end{equation}

\section{Social Media Sentiment Analysis}

\subsection{Sentiment Data Schema}

Social media data provides real-time insights into market sentiment and community engagement.

\begin{lstlisting}
{
    "sentiment_data": {
        "post_id": "1234567890123456789",
        "platform": "twitter",
        "text": "Just bought more $ETH! Bullish on DeFi 🚀",
        "timestamp": "2024-01-15T10:30:00Z",
        "user_id": "987654321",
        "user_followers": 50000,
        "user_verified": true,
        "sentiment_score": 0.85,
        "engagement_score": 0.72,
        "hashtags": ["#ETH", "#DeFi", "#crypto"],
        "mentions": ["@VitalikButerin"],
        "retweets": 150,
        "likes": 500,
        "replies": 25
    }
}
\end{lstlisting}

\subsection{Sentiment Scoring Methods}

\subsubsection{VADER Sentiment Analysis}

VADER (Valence Aware Dictionary and sEntiment Reasoner) is specifically attuned to sentiments expressed in social media.

\begin{equation}
\text{VADER Score} = \frac{\text{Positive} - \text{Negative}}{\text{Positive} + \text{Negative} + \text{Neutral}}
\end{equation}

\subsubsection{Compound Sentiment Score}

The compound score is a normalized score between -1 (most negative) and +1 (most positive).

\begin{equation}
\text{Compound Score} = \frac{\text{Sum of VADER Scores}}{\text{Number of Words}}
\end{equation}

\subsection{Engagement Analysis}

Engagement analysis measures the virality and impact of social media content.

\begin{equation}
\text{Engagement Rate} = \frac{\text{Likes} + \text{Retweets} + \text{Replies}}{\text{Followers}} \times 100\%
\end{equation}

\begin{equation}
\text{Viral Coefficient} = \frac{\text{Retweets} \times \text{Average Followers of Retweeters}}{\text{Followers}}
\end{equation}

\begin{equation}
\text{Influence Score} = \log(\text{Followers}) \times \text{Engagement Rate} \times \text{Verified Factor}
\end{equation}

\subsection{Trending Topic Analysis}

Trending topic analysis identifies emerging themes and discussions.

\begin{equation}
\text{Topic Frequency} = \frac{\text{Mentions of Topic}_{24h}}{\text{Total Posts}_{24h}} \times 100\%
\end{equation}

\begin{equation}
\text{Topic Momentum} = \frac{\text{Topic Frequency}_{t} - \text{Topic Frequency}_{t-24h}}{\text{Topic Frequency}_{t-24h}}
\end{equation}

\section{News Media Analysis}

\subsection{News Data Schema}

News media data provides insights into market sentiment and regulatory developments.

\begin{lstlisting}
{
    "news_data": {
        "article_id": "12345",
        "title": "Ethereum Layer 2 Solutions See Record Growth",
        "content": "The total value locked in Ethereum Layer 2...",
        "source": "CoinDesk",
        "published_at": "2024-01-15T10:30:00Z",
        "sentiment_score": 0.75,
        "impact_score": 0.82,
        "keywords": ["ethereum", "layer2", "defi", "growth"],
        "read_time": 300,
        "word_count": 1200,
        "category": "technology"
    }
}
\end{lstlisting}

\subsection{News Sentiment Analysis}

News sentiment analysis examines the tone and impact of media coverage.

\begin{equation}
\text{News Sentiment Score} = \frac{\text{Positive Articles} - \text{Negative Articles}}{\text{Total Articles}}
\end{equation}

\begin{equation}
\text{News Momentum} = \frac{\text{Sentiment}_{t} - \text{Sentiment}_{t-24h}}{\text{Sentiment}_{t-24h}}
\end{equation}

\subsection{Impact Analysis}

Impact analysis measures the potential market impact of news events.

\begin{equation}
\text{Impact Score} = \text{Sentiment Score} \times \text{Source Credibility} \times \text{Read Time}
\end{equation}

\begin{equation}
\text{News Volume} = \frac{\text{Articles}_{24h}}{\text{Total Articles}_{7d}} \times 100\%
\end{equation}

\section{Feature Engineering for Machine Learning}

\subsection{Time Series Feature Engineering}

\subsubsection{Lag Features}

Lag features capture temporal dependencies in the data.

\begin{equation}
\text{Lag Feature}_k = X_{t-k}
\end{equation}

\begin{lstlisting}
def create_lag_features(data, columns, lags=[1, 6, 12, 24]):
    """Create lag features for time series data"""
    for col in columns:
        for lag in lags:
            data[f'{col}_lag_{lag}'] = data[col].shift(lag)
    return data
\end{lstlisting}

\subsubsection{Rolling Window Features}

Rolling window features capture local trends and patterns.

\begin{equation}
\text{Rolling Mean}_k = \frac{1}{k} \sum_{i=t-k+1}^{t} X_i
\end{equation}

\begin{equation}
\text{Rolling Std}_k = \sqrt{\frac{1}{k-1} \sum_{i=t-k+1}^{t} (X_i - \bar{X})^2}
\end{equation}

\subsubsection{Exponential Moving Averages}

Exponential moving averages provide weighted averages with more emphasis on recent data.

\begin{equation}
\text{EMA}_t = \alpha \times X_t + (1-\alpha) \times \text{EMA}_{t-1}
\end{equation}

where $\alpha = \frac{2}{k+1}$ and $k$ is the window size.

\subsection{Cross-Source Feature Engineering}

\subsubsection{Correlation Features}

Correlation features capture relationships between different data sources.

\begin{equation}
\text{Price-Sentiment Correlation} = \text{Corr}(\text{Price}_{24h}, \text{Sentiment}_{24h})
\end{equation}

\begin{equation}
\text{Volume-Development Correlation} = \text{Corr}(\text{Volume}_{24h}, \text{Commit Frequency}_{7d})
\end{equation}

\subsubsection{Interaction Features}

Interaction features capture non-linear relationships between variables.

\begin{equation}
\text{Sentiment Volume Interaction} = \text{Sentiment Score} \times \text{Volume Change}
\end{equation}

\begin{equation}
\text{Development Momentum} = \text{Commit Frequency} \times \text{Star Growth Rate}
\end{equation}

\section{Statistical Analysis Methods}

\subsection{Descriptive Statistics}

\subsubsection{Central Tendency Measures}

\begin{equation}
\text{Mean} = \bar{X} = \frac{1}{n} \sum_{i=1}^{n} X_i
\end{equation}

\begin{equation}
\text{Median} = \text{Value at } \frac{n+1}{2} \text{ position}
\end{equation}

\begin{equation}
\text{Mode} = \text{Most frequent value}
\end{equation}

\subsubsection{Dispersion Measures}

\begin{equation}
\text{Variance} = \sigma^2 = \frac{1}{n-1} \sum_{i=1}^{n} (X_i - \bar{X})^2
\end{equation}

\begin{equation}
\text{Standard Deviation} = \sigma = \sqrt{\sigma^2}
\end{equation}

\begin{equation}
\text{Coefficient of Variation} = \text{CV} = \frac{\sigma}{\bar{X}} \times 100\%
\end{equation}

\subsection{Correlation Analysis}

\subsubsection{Pearson Correlation}

\begin{equation}
r = \frac{\sum_{i=1}^{n} (X_i - \bar{X})(Y_i - \bar{Y})}{\sqrt{\sum_{i=1}^{n} (X_i - \bar{X})^2} \sqrt{\sum_{i=1}^{n} (Y_i - \bar{Y})^2}}
\end{equation}

\subsubsection{Spearman Rank Correlation}

\begin{equation}
\rho = 1 - \frac{6 \sum d_i^2}{n(n^2-1)}
\end{equation}

where $d_i$ is the difference between ranks.

\subsection{Time Series Analysis}

\subsubsection{Stationarity Tests}

Augmented Dickey-Fuller test for stationarity:

\begin{equation}
\Delta Y_t = \alpha + \beta t + \gamma Y_{t-1} + \sum_{i=1}^{p} \delta_i \Delta Y_{t-i} + \epsilon_t
\end{equation}

\subsubsection{Autocorrelation Analysis}

\begin{equation}
\text{ACF}(k) = \frac{\text{Cov}(Y_t, Y_{t-k})}{\text{Var}(Y_t)}
\end{equation}

\section{Data Quality Assessment}

\subsection{Completeness Analysis}

\begin{equation}
\text{Completeness Rate} = \frac{\text{Non-null Values}}{\text{Total Values}} \times 100\%
\end{equation}

\subsection{Accuracy Assessment}

\begin{equation}
\text{Accuracy Score} = \frac{\text{Correct Predictions}}{\text{Total Predictions}} \times 100\%
\end{equation}

\subsection{Consistency Checks}

\begin{equation}
\text{Consistency Score} = \frac{\text{Consistent Records}}{\text{Total Records}} \times 100\%
\end{equation}

\subsection{Freshness Analysis}

\begin{equation}
\text{Freshness Score} = \frac{\text{Recent Data Points}}{\text{Total Data Points}} \times 100\%
\end{equation}

\section{Machine Learning Preparation}

\subsection{Dataset Construction}

\subsubsection{Feature Matrix Construction}

The feature matrix $X$ is constructed as:

\begin{equation}
X = \begin{bmatrix}
x_{1,1} & x_{1,2} & \cdots & x_{1,p} \\
x_{2,1} & x_{2,2} & \cdots & x_{2,p} \\
\vdots & \vdots & \ddots & \vdots \\
x_{n,1} & x_{n,2} & \cdots & x_{n,p}
\end{bmatrix}
\end{equation}

where $n$ is the number of samples and $p$ is the number of features.

\subsubsection{Target Variable Construction}

Target variables are constructed for different prediction tasks:

\begin{equation}
y_{price} = \text{Price}_{t+h} - \text{Price}_t
\end{equation}

\begin{equation}
y_{risk} = \text{Risk Category} \in \{0, 1, 2, 3, 4\}
\end{equation}

\begin{equation}
y_{sentiment} = \text{Sentiment Class} \in \{0, 1, 2\}
\end{equation}

\subsection{Data Preprocessing}

\subsubsection{Standardization}

\begin{equation}
X_{standardized} = \frac{X - \mu}{\sigma}
\end{equation}

\subsubsection{Normalization}

\begin{equation}
X_{normalized} = \frac{X - X_{min}}{X_{max} - X_{min}}
\end{equation}

\subsubsection{Missing Value Imputation}

\begin{equation}
X_{imputed} = \text{Mean}(X_{non-null}) \text{ for missing values}
\end{equation}

\section{Visualization Techniques}

\subsection{Time Series Visualization}

\subsubsection{Price and Volume Charts}

\begin{lstlisting}
def plot_price_volume(data):
    """Plot price and volume time series"""
    fig, (ax1, ax2) = plt.subplots(2, 1, figsize=(15, 10))
    
    # Price chart
    ax1.plot(data['timestamp'], data['price'], 
             color='blue', linewidth=1.5)
    ax1.set_title('Price Trend (24h)', fontsize=14, fontweight='bold')
    ax1.set_ylabel('Price (USD)', fontsize=12)
    ax1.grid(True, alpha=0.3)
    
    # Volume chart
    ax2.bar(data['timestamp'], data['volume'], 
            color='green', alpha=0.7)
    ax2.set_title('Volume Trend (24h)', fontsize=14, fontweight='bold')
    ax2.set_ylabel('Volume (USD)', fontsize=12)
    ax2.set_xlabel('Time', fontsize=12)
    ax2.grid(True, alpha=0.3)
    
    plt.tight_layout()
    plt.show()
\end{lstlisting}

\subsection{Correlation Analysis Visualization}

\subsubsection{Correlation Heatmap}

\begin{lstlisting}
def plot_correlation_heatmap(data, features):
    """Plot correlation heatmap"""
    correlation_matrix = data[features].corr()
    
    plt.figure(figsize=(20, 16))
    sns.heatmap(correlation_matrix, 
                annot=True, 
                cmap='coolwarm', 
                center=0,
                square=True,
                fmt='.2f',
                cbar_kws={'shrink': 0.8})
    
    plt.title('Feature Correlation Heatmap', 
              fontsize=16, fontweight='bold')
    plt.tight_layout()
    plt.show()
\end{lstlisting}

\subsection{Distribution Analysis}

\subsubsection{Histogram and Density Plots}

\begin{lstlisting}
def plot_distribution_analysis(data, features):
    """Plot distribution analysis"""
    fig, axes = plt.subplots(2, 2, figsize=(15, 12))
    axes = axes.ravel()
    
    for i, feature in enumerate(features):
        # Histogram
        axes[i].hist(data[feature], bins=50, alpha=0.7, 
                     color='skyblue', edgecolor='black')
        axes[i].set_title(f'{feature} Distribution', 
                          fontsize=12, fontweight='bold')
        axes[i].set_xlabel(feature, fontsize=10)
        axes[i].set_ylabel('Frequency', fontsize=10)
        axes[i].grid(True, alpha=0.3)
    
    plt.tight_layout()
    plt.show()
\end{lstlisting}

\section{Performance Metrics}

\subsection{Regression Metrics}

\subsubsection{Mean Absolute Error (MAE)}

\begin{equation}
\text{MAE} = \frac{1}{n} \sum_{i=1}^{n} |y_i - \hat{y}_i|
\end{equation}

\subsubsection{Mean Squared Error (MSE)}

\begin{equation}
\text{MSE} = \frac{1}{n} \sum_{i=1}^{n} (y_i - \hat{y}_i)^2
\end{equation}

\subsubsection{Root Mean Squared Error (RMSE)}

\begin{equation}
\text{RMSE} = \sqrt{\frac{1}{n} \sum_{i=1}^{n} (y_i - \hat{y}_i)^2}
\end{equation}

\subsubsection{Mean Absolute Percentage Error (MAPE)}

\begin{equation}
\text{MAPE} = \frac{100\%}{n} \sum_{i=1}^{n} \left|\frac{y_i - \hat{y}_i}{y_i}\right|
\end{equation}

\subsection{Classification Metrics}

\subsubsection{Accuracy}

\begin{equation}
\text{Accuracy} = \frac{\text{TP} + \text{TN}}{\text{TP} + \text{TN} + \text{FP} + \text{FN}}
\end{equation}

\subsubsection{Precision}

\begin{equation}
\text{Precision} = \frac{\text{TP}}{\text{TP} + \text{FP}}
\end{equation}

\subsubsection{Recall}

\begin{equation}
\text{Recall} = \frac{\text{TP}}{\text{TP} + \text{FN}}
\end{equation}

\subsubsection{F1-Score}

\begin{equation}
\text{F1-Score} = 2 \times \frac{\text{Precision} \times \text{Recall}}{\text{Precision} + \text{Recall}}
\end{equation}

\section{Conclusion}

This comprehensive data analytics guide provides a framework for analyzing multi-source blockchain data for machine learning applications. The guide covers data collection, preprocessing, feature engineering, and statistical analysis methods specifically designed for cryptocurrency and DeFi analytics.

Key takeaways for data analysts:

\begin{enumerate}
    \item \textbf{Data Quality}: Maintain 93.3\% overall data completeness across all sources
    \item \textbf{Feature Engineering}: Create 125 engineered features from 5 data sources
    \item \textbf{Time Series Analysis}: Implement lag features and rolling windows for temporal patterns
    \item \textbf{Cross-Source Integration}: Leverage correlations between different data types
    \item \textbf{Performance Monitoring}: Track accuracy, precision, and recall for model evaluation
\end{enumerate}

The framework enables data analysts to build robust machine learning models for price prediction, risk assessment, and sentiment analysis in the blockchain and DeFi space.

\end{document}
